\documentclass[12pt,a4paper]{article}

\usepackage{amsmath,amssymb}
\usepackage{geometry}
\usepackage{graphicx}
\usepackage{setspace}
\geometry{margin=1in}
\setstretch{1.2}

\title{CSE400 -- Fundamentals of Probability in Computing\\
Lecture Scribe}
\author{Joint Random Variables, Joint Distributions, and Joint CDF\\
Source: Lecture L15--L18}
\date{}

\begin{document}

\maketitle

\section{Introduction: Why Study Joint Random Variables?}

\subsection{Core Idea}

In many real-world problems we are interested in two or more random variables simultaneously. Instead of studying variables separately, we analyze their relationship.

Example: Height of a person and weight of a person.  
Both variables vary randomly but are also related.

\subsection{Examples from the Lecture}

\textbf{Example 1: Giraffes}

Random variables:
\[
X = \text{Height of giraffe}, \quad Y = \text{Weight of giraffe}
\]

Observation: Taller giraffes tend to weigh more.

\textbf{Example 2: Children Data}

\[
X = \text{IQ}, \quad Y = \text{Birth weight}
\]

Researchers study how early health conditions influence intelligence.

\textbf{Example 3: Exercise and Heart Disease}

\[
X = \text{Frequency of exercise}, \quad Y = \text{Heart disease rate}
\]

More exercise usually leads to lower heart disease risk.

\textbf{Example 4: Air Pollution}

\[
X = \text{Pollution level}, \quad Y = \text{Respiratory illness rate}
\]

Higher pollution levels often increase illness rates.

\textbf{Example 5: Social Media}

\[
X = \text{Number of Facebook friends}, \quad Y = \text{Age of user}
\]

Different age groups exhibit different online behavior.

\subsection{Engineering Applications}

\textbf{Smart Transportation}

\[
X = \text{Vehicle speed}, \quad Y = \text{Traffic density}
\]

Higher traffic density usually leads to lower vehicle speed.

Applications include adaptive traffic signals, congestion prediction, and autonomous vehicle routing.

\textbf{Wireless Communication}

\[
X = \text{Signal-to-noise ratio (SNR)}, \quad Y = \text{Data throughput}
\]

Higher SNR generally produces higher data rates.

Applications include 5G scheduling, adaptive modulation, and network planning.

\textbf{Weather Prediction}

\[
X = \text{Rainfall}, \quad Y = \text{Temperature}
\]

Meteorologists study relationships between weather variables.

\textbf{Drone Delivery}

\[
X = \text{Wind speed}, \quad Y = \text{Delivery time}
\]

Higher wind speeds increase delivery times.

Applications include UAV logistics and autonomous delivery systems.

\subsection{Why Joint Distributions Are Important}

Joint distributions help us understand:

\begin{itemize}
\item Probability of two events occurring together
\item Relationships between variables
\item Dependence or independence between variables
\end{itemize}

To measure relationships we use covariance and correlation.

\section{Discrete Case: Joint PMF}

\subsection{Setup}

Suppose we have two discrete random variables:

\[
X = \{x_1, x_2, \ldots, x_n\}
\]

\[
Y = \{y_1, y_2, \ldots, y_m\}
\]

The pair $(X,Y)$ can take values:

\[
(x_1,y_1),(x_1,y_2),\ldots,(x_n,y_m)
\]

This set of all combinations is called the \textbf{product space}.

\subsection{Definition: Joint PMF}

The joint probability mass function is defined as

\[
p(x_i,y_j)=P(X=x_i,Y=y_j)
\]

This represents the probability that both events occur simultaneously.

\subsection{Joint PMF Table Representation}

\[
\begin{array}{c|cccc}
X\backslash Y & y_1 & y_2 & \cdots & y_m\\
\hline
x_1 & p(x_1,y_1) & p(x_1,y_2) & \cdots & p(x_1,y_m)\\
x_2 & p(x_2,y_1) & p(x_2,y_2) & \cdots & p(x_2,y_m)\\
\vdots & \vdots & \vdots & \ddots & \vdots\\
x_n & p(x_n,y_1) & p(x_n,y_2) & \cdots & p(x_n,y_m)
\end{array}
\]

\subsection{Properties of Joint PMF}

\textbf{Property 1}

\[
0 \le p(x_i,y_j) \le 1
\]

\textbf{Property 2}

\[
\sum_i \sum_j p(x_i,y_j) = 1
\]

\section{Example: Rolling Two Dice}

Experiment: Roll two fair dice.

Define random variables:

\[
X = \text{value on the first die}
\]

\[
T = \text{total sum of both dice}
\]

Each outcome $(i,j)$ has probability:

\[
\frac{1}{36}
\]

since there are $36$ possible outcomes.

Example probabilities:

\[
P(X=1,T=2)=\frac{1}{36}
\]

\[
P(X=1,T=3)=\frac{1}{36}
\]

\[
P(X=1,T=8)=0
\]

because the sum cannot equal $8$ when the first die is $1$.

Observation: $X$ and $T$ are not independent because $T$ depends on $X$.

\section{Continuous Case: Joint PDF}

For continuous variables, the joint distribution is described using a \textbf{joint probability density function (PDF)}.

Suppose

\[
X \in [a,b], \quad Y \in [c,d]
\]

\subsection{Definition}

A function $f(x,y)$ is called a joint PDF if

\[
P((x,y)\text{ in small region}) \approx f(x,y)\,dx\,dy
\]

where $dx\,dy$ represents the area of a small rectangle.

\subsection{Properties}

\textbf{Property 1}

\[
f(x,y) \ge 0
\]

\textbf{Property 2}

\[
\int\int f(x,y)\,dx\,dy = 1
\]

over the entire region.

Note: The value of $f(x,y)$ can be greater than $1$ because it represents density, not probability.

\section{Worked Example}

Given:

\[
f(x,y)=4xy
\]

for

\[
0 \le x \le 1, \quad 0 \le y \le 1
\]

\subsection{Step 1: Check Normalization}

\[
\int_0^1 \int_0^1 4xy \, dx\,dy
\]

Inner integral:

\[
\int_0^1 4xy\,dx = 4y\int_0^1 x\,dx
\]

\[
=4y \cdot \frac{1}{2}=2y
\]

Outer integral:

\[
\int_0^1 2y\,dy = 2\cdot\frac{1}{2}=1
\]

Thus $f(x,y)$ is a valid joint PDF.

\subsection{Step 2: Compute Event Probability}

Event:

\[
A = \{X<0.5,\;Y>0.5\}
\]

\[
P(A)=\int_0^{0.5}\int_{0.5}^{1}4xy\,dy\,dx
\]

Inner integral:

\[
\int_{0.5}^{1}4xy\,dy
\]

\[
=4x\left[\frac{y^2}{2}\right]_{0.5}^{1}
\]

\[
=2x(1-0.25)=1.5x
\]

Outer integral:

\[
\int_0^{0.5}1.5x\,dx
\]

\[
=1.5\left[\frac{x^2}{2}\right]_0^{0.5}
\]

\[
=\frac{3}{16}
\]

\[
P(A)=\frac{3}{16}
\]

\section{Joint Cumulative Distribution Function}

For two random variables, the joint CDF is defined as

\[
F(x,y)=P(X\le x,Y\le y)
\]

\subsection{Continuous Case}

If the joint PDF is known:

\[
F(x,y)=\int_{-\infty}^{x}\int_{-\infty}^{y}f(u,v)\,du\,dv
\]

\subsection{Recovering Joint PDF}

\[
f(x,y)=\frac{\partial^2 F(x,y)}{\partial x \partial y}
\]

\subsection{Discrete Case}

For discrete variables:

\[
F(x,y)=\sum_{x_i\le x}\sum_{y_j\le y}p(x_i,y_j)
\]

This means adding all probabilities in the table satisfying those conditions.

\section{Properties of Joint CDF}

\textbf{Property 1: Non-decreasing}

$F(x,y)$ never decreases as $x$ or $y$ increases.

\textbf{Property 2: Lower Left Corner}

\[
F(x,y)=0
\]

when

\[
x\to -\infty, \quad y\to -\infty
\]

\textbf{Property 3: Upper Right Corner}

\[
F(x,y)=1
\]

when

\[
x\to \infty, \quad y\to \infty
\]

Conceptually, the joint CDF represents the area under the joint PDF surface up to the point $(x,y)$.

\section{Summary}

Key concepts covered:

\begin{itemize}
\item Joint random variables
\item Joint PMF
\item Joint PDF
\item Joint probability tables
\item Joint cumulative distribution function
\item Continuous vs discrete joint distributions
\item Probability computation using integrals
\item Engineering applications
\end{itemize}

\section{Formula Sheet}

\subsection*{Joint PMF}

\[
p(x,y)=P(X=x,Y=y)
\]

\[
0 \le p(x,y) \le 1
\]

\[
\sum\sum p(x,y)=1
\]

\subsection*{Joint PDF}

\[
f(x,y)\ge0
\]

\[
\int\int f(x,y)\,dx\,dy=1
\]

\subsection*{Probability of Region}

\[
P(A)=\int\int_A f(x,y)\,dx\,dy
\]

\subsection*{Joint CDF}

\[
F(x,y)=P(X\le x,Y\le y)
\]

Continuous case:

\[
F(x,y)=\int_{-\infty}^{x}\int_{-\infty}^{y}f(u,v)\,du\,dv
\]

Recover PDF:

\[
f(x,y)=\frac{\partial^2 F(x,y)}{\partial x \partial y}
\]

Discrete case:

\[
F(x,y)=\sum_{x_i\le x}\sum_{y_j\le y}p(x_i,y_j)
\]

\end{document}